\section{The whole selected chart and the pointwise source identity}
\label{sec:v70-whole-chart}
We now apply the overlap comparison to the entire selected Morse
disk, not only to an inner germ collar. This accounts for its outer
annuli on the same footing as its center. The genuinely uncovered
physical source remains separate.

For each fixed $\chi,\varepsilon$ choose
\[
 I_\chi=(0,h_\chi),\qquad J_\chi=(0,h_\chi/2),
 \qquad
 s_{70}=\sum_zs_z^{I_\chi,J_\chi},\quad
 d_{70}=\sum_zd_z^{I_\chi,J_\chi}.
\]
The sums are over all selected $\chi$-incident physical charts with
the prescribed exact label. Write
$e_{70}=b^{\varepsilon,\chi,\mathrm{outside}}$ for the original
positive complement in \eqref{eq:v66-complete-positive-partition}.

\begin{theorem}[Complete allocation of the selected-chart source]
\label{thm:v70-whole-chart-allocation}
On the original probability space, and on the corresponding common
coarea versions,
\begin{equation}\label{eq:v70-complete-source-identity}
 b^{\varepsilon,\mathrm{clr}}=s_{70}+d_{70}+e_{70},
 \qquad s_{70},d_{70},e_{70}\ge0.
\end{equation}
The selected-chart and outside events are disjoint; within each chart
$s_{70}$ and $d_{70}$ are complementary positive weights. Moreover
\begin{equation}\label{eq:v70-whole-chart-height}
 \mathcal H(s_{70})
 \le\frac32Ch_\chi
       \exp\{(1+1/\sqrt2)K_\chi r_\chi\}
 \le C'\chi^6.
\end{equation}
The constants are independent of the exact label and the collision
count. Every portion of every selected disk, including all positive-
radius births and all its outer annuli, is allocated in
\eqref{eq:v70-complete-source-identity}. Neither $d_{70}$ nor
$e_{70}$ is assigned a small height by this assertion.
\end{theorem}
\begin{proof}
Equation~\eqref{eq:v70-local-source-split}, summed over the disjoint
physical chart restrictions, is precisely the chart source. Add the
unchanged outside source to obtain \eqref{eq:v70-complete-source-identity}.
Apply Theorem~\ref{thm:v70-overlap-height} with
$(a,b,c,d)=(0,h_\chi,0,h_\chi/2)$ and $\delta=1$. The ratio
$d/(d-c)$ is one, the center interval plus the collar width has
length $3h_\chi/2$, and the distortion exponent is
$(1+1/\sqrt2)K_\chi r_\chi$. The source interval is open, so
$b=h_\chi$ only uses estimates strictly within the Morse disk.
The receiving width is strictly smaller than $h_\chi$, as required
by the inherited collar theorem. Finally
$h_\chi=r_0^2\chi^6/2$ and $K_\chi r_\chi=Cr_0\chi$.
This proves the uniform last bound. No exhaustion in the collision
count or interchange with an infinite sum of annuli occurs.
\end{proof}

The complement $e_{70}$ retains failed taper, selected grazing,
noncritical roof-rank pieces, and every state outside the specified
physical disks. Unlike $e_{69}$ it does not contain annuli cut out
of these disks at an additional small radial cutoff. Those annuli
are instead divided, with their original weight, between $s_{70}$
and $d_{70}$. They have not all been declared small: a late birth
with no occupation in $J_\chi$ is allocated wholly to $d_{70}$.
This example also explains why the complete radial tail is not proved
merely by introducing the overlap decomposition.

\begin{corollary}[An ordered whole-chart contribution to raw inversion]
\label{cor:v70-ordered-raw-budget}
For every sufficiently large fixed reconstruction band $B$, put
\[
 \varepsilon(B)=A_0B^{-1/12},\qquad
                  \chi(B)=\sqrt{\varepsilon(B)}.
\]
These choices are admissible for all sufficiently large $B$.
The unchanged central pointwise error satisfies
\begin{equation}\label{eq:v70-ordered-raw-budget}
 \begin{split}
 \mathcal E_M\le{}&C_MB^{-1/192}+C'B^{-1/4}
           +\mathcal H(b^{\varepsilon(B),\mathrm{inc}})\\
           &+\mathcal H(d_{70})+\mathcal H(e_{70}).
 \end{split}
\end{equation}
Here $s_{70},d_{70},e_{70}$ are formed at the fixed choices
$\chi(B),\varepsilon(B)$ before taking the collision limsup.
No cutoff in this statement depends on $m$.
\end{corollary}
\begin{proof}
For sufficiently small $\varepsilon$, one has
$4\varepsilon\le\sqrt\varepsilon\le1/10$. Thus all chart and
incidence assumptions hold. Theorem~\ref{thm:v51-positive-error},
the original first-incidence/first-clearance split, and
\eqref{eq:v70-complete-source-identity} give the asserted inequality
with $\mathcal H(s_{70})$ in place of $C'B^{-1/4}$.
By \eqref{eq:v70-whole-chart-height}, that height is at most
$C'\chi(B)^6=C'\varepsilon(B)^3=C'A_0^3B^{-1/4}$.
Absorb the fixed $A_0^3$ into the constant. Only after this fixed-band
collision limit may $B$ tend to infinity.
\end{proof}

The unrestricted two-sided arithmetic theorem would follow from
vanishing of the last three heights in
\eqref{eq:v70-ordered-raw-budget}. These estimates are not consequences
of finite-word BV, of the radial-moment implication, or of the
finite-count bound $CA^m\varepsilon$ for first incidence. In
particular no assertion of a finite uniform radial moment or a
vanishing full signed-current tail is made here. The scalar endpoint,
its exact arithmetic zero classes, and the original source remain
unchanged. Unrestricted same-roof bridges, forward essential likelihood
and pointwise roof-conditioned path laws still require that endpoint.

\subsection{Comparison with existing local-limit statements}
The local Lorentz-process theorem of \cite{SV} already gives a
cell-index local limit. Theorem 2.2 and Remark 2.3 of \cite{DPZ}
give a two-dimensional cell-index local limit with admissible endpoint
observables, including piecewise H\"older endpoints under the stated
uniformity conditions. Inserting a free-flight observable as an
endpoint factor in that theorem is different from simultaneously
resolving the accumulated roof, exact section occupation and the
present density-level boundary source. We do not infer the latter
from an endpoint insertion.

Definition 2.2 of \cite{DN} formulates its mixing local limit using
compactly supported continuous tests in the additive coordinate;
Remark 2.5 permits appropriate continuity-set indicators. Its
Proposition 6.3 applies the suspension theory to finite-horizon Sinai
billiards under the joint minimality and arithmetic-group hypotheses.
Those statements supply neither an essential-supremum estimate for
$\Delta_{z,J}$ nor the complete three-height estimate in
\eqref{eq:v70-ordered-raw-budget}. Differentiating an interval law
would require precisely an additional density-height argument.
The present overlap and signed-current results address this source
question; they are not claims of a new general suspension mechanism
or a replacement for the inherited anisotropic local-limit inputs.
\label{end:v70-new-mathematics}
